% Options for packages loaded elsewhere
\PassOptionsToPackage{unicode}{hyperref}
\PassOptionsToPackage{hyphens}{url}
\documentclass[
  journal]{IEEEtran}
\usepackage{xcolor}
\usepackage{amsmath,amssymb}
\usepackage{iftex}
\ifPDFTeX
  \usepackage[T1]{fontenc}
  \usepackage[utf8]{inputenc}
  \usepackage{textcomp} % provide euro and other symbols
\else % if luatex or xetex
  \usepackage{unicode-math} % this also loads fontspec
  \defaultfontfeatures{Scale=MatchLowercase}
  \defaultfontfeatures[\rmfamily]{Ligatures=TeX,Scale=1}
\fi
\usepackage{lmodern}
\ifPDFTeX\else
  % xetex/luatex font selection
  \setmainfont[]{Times New Roman}
\fi
% Use upquote if available, for straight quotes in verbatim environments
\IfFileExists{upquote.sty}{\usepackage{upquote}}{}
\IfFileExists{microtype.sty}{% use microtype if available
  \usepackage[]{microtype}
  \UseMicrotypeSet[protrusion]{basicmath} % disable protrusion for tt fonts
}{}
\makeatletter
\@ifundefined{KOMAClassName}{% if non-KOMA class
  \IfFileExists{parskip.sty}{%
    \usepackage{parskip}
  }{% else
    \setlength{\parindent}{0pt}
    \setlength{\parskip}{6pt plus 2pt minus 1pt}}
}{% if KOMA class
  \KOMAoptions{parskip=half}}
\makeatother
\usepackage{color}
\usepackage{fancyvrb}
\newcommand{\VerbBar}{|}
\newcommand{\VERB}{\Verb[commandchars=\\\{\}]}
\DefineVerbatimEnvironment{Highlighting}{Verbatim}{commandchars=\\\{\}}
% Add ',fontsize=\small' for more characters per line
\newenvironment{Shaded}{}{}
\newcommand{\AlertTok}[1]{\textcolor[rgb]{1.00,0.00,0.00}{\textbf{#1}}}
\newcommand{\AnnotationTok}[1]{\textcolor[rgb]{0.38,0.63,0.69}{\textbf{\textit{#1}}}}
\newcommand{\AttributeTok}[1]{\textcolor[rgb]{0.49,0.56,0.16}{#1}}
\newcommand{\BaseNTok}[1]{\textcolor[rgb]{0.25,0.63,0.44}{#1}}
\newcommand{\BuiltInTok}[1]{\textcolor[rgb]{0.00,0.50,0.00}{#1}}
\newcommand{\CharTok}[1]{\textcolor[rgb]{0.25,0.44,0.63}{#1}}
\newcommand{\CommentTok}[1]{\textcolor[rgb]{0.38,0.63,0.69}{\textit{#1}}}
\newcommand{\CommentVarTok}[1]{\textcolor[rgb]{0.38,0.63,0.69}{\textbf{\textit{#1}}}}
\newcommand{\ConstantTok}[1]{\textcolor[rgb]{0.53,0.00,0.00}{#1}}
\newcommand{\ControlFlowTok}[1]{\textcolor[rgb]{0.00,0.44,0.13}{\textbf{#1}}}
\newcommand{\DataTypeTok}[1]{\textcolor[rgb]{0.56,0.13,0.00}{#1}}
\newcommand{\DecValTok}[1]{\textcolor[rgb]{0.25,0.63,0.44}{#1}}
\newcommand{\DocumentationTok}[1]{\textcolor[rgb]{0.73,0.13,0.13}{\textit{#1}}}
\newcommand{\ErrorTok}[1]{\textcolor[rgb]{1.00,0.00,0.00}{\textbf{#1}}}
\newcommand{\ExtensionTok}[1]{#1}
\newcommand{\FloatTok}[1]{\textcolor[rgb]{0.25,0.63,0.44}{#1}}
\newcommand{\FunctionTok}[1]{\textcolor[rgb]{0.02,0.16,0.49}{#1}}
\newcommand{\ImportTok}[1]{\textcolor[rgb]{0.00,0.50,0.00}{\textbf{#1}}}
\newcommand{\InformationTok}[1]{\textcolor[rgb]{0.38,0.63,0.69}{\textbf{\textit{#1}}}}
\newcommand{\KeywordTok}[1]{\textcolor[rgb]{0.00,0.44,0.13}{\textbf{#1}}}
\newcommand{\NormalTok}[1]{#1}
\newcommand{\OperatorTok}[1]{\textcolor[rgb]{0.40,0.40,0.40}{#1}}
\newcommand{\OtherTok}[1]{\textcolor[rgb]{0.00,0.44,0.13}{#1}}
\newcommand{\PreprocessorTok}[1]{\textcolor[rgb]{0.74,0.48,0.00}{#1}}
\newcommand{\RegionMarkerTok}[1]{#1}
\newcommand{\SpecialCharTok}[1]{\textcolor[rgb]{0.25,0.44,0.63}{#1}}
\newcommand{\SpecialStringTok}[1]{\textcolor[rgb]{0.73,0.40,0.53}{#1}}
\newcommand{\StringTok}[1]{\textcolor[rgb]{0.25,0.44,0.63}{#1}}
\newcommand{\VariableTok}[1]{\textcolor[rgb]{0.10,0.09,0.49}{#1}}
\newcommand{\VerbatimStringTok}[1]{\textcolor[rgb]{0.25,0.44,0.63}{#1}}
\newcommand{\WarningTok}[1]{\textcolor[rgb]{0.38,0.63,0.69}{\textbf{\textit{#1}}}}
\usepackage{longtable,booktabs,array}
\usepackage{calc} % for calculating minipage widths
% Correct order of tables after \paragraph or \subparagraph
\usepackage{etoolbox}
\makeatletter
\patchcmd\longtable{\par}{\if@noskipsec\mbox{}\fi\par}{}{}
\makeatother
% Allow footnotes in longtable head/foot
\IfFileExists{footnotehyper.sty}{\usepackage{footnotehyper}}{\usepackage{footnote}}
\makesavenoteenv{longtable}
\setlength{\emergencystretch}{3em} % prevent overfull lines
\providecommand{\tightlist}{%
  \setlength{\itemsep}{0pt}\setlength{\parskip}{0pt}}
\usepackage{bookmark}
\IfFileExists{xurl.sty}{\usepackage{xurl}}{} % add URL line breaks if available
\urlstyle{same}
\hypersetup{
  hidelinks,
  pdfcreator={LaTeX via pandoc}}

\author{}
\date{}

\begin{document}

\section{Natural Language as a Tactical Control Interface for Real-Time
Strategy Games: An Engine-FSM Minimal MCP Approach Evaluated on
OpenRA}\label{natural-language-as-a-tactical-control-interface-for-real-time-strategy-games-an-engine-fsm-minimal-mcp-approach-evaluated-on-openra}

JiZiYi\\
Universiti Teknologi Malaysia (UTM)\\
jiziyi@graduate.utm.my

\subsection{Abstract}\label{abstract}

Real-time strategy (RTS) games coordinate dozens of mobile units across
simultaneous fronts. Even when the player's tactical decision is clear,
executing it through selection boxes, control groups, and queued orders
is mechanically costly. We ask whether plain natural language can serve
as the tactical control interface: can a player say ``send the tanks
down the middle and have the APCs flank from both sides'' and have 30
mixed units execute the maneuver, without naming actor IDs, computing
coordinates, or selecting from a 50-tool API menu? We present
\texttt{openra\_mcp}, a co-pilot for OpenRA --- a production RTS engine
--- that exposes only two engine-side execution primitives
(\texttt{Assault}, \texttt{Protection}) over the Model Context Protocol
(MCP). The player issues free-form natural-language commands; an LLM
translates them into structured squad calls; the engine runs per-tick
movement, pathfinding, and combat. Higher-level tactics --- pincer,
multi-squad split, feint-and-raid, route constraints, time-sequenced
attacks --- are composed by the LLM on top of the two primitives rather
than encoded as specialized engine FSMs. We evaluate the interface on
three axes: (i) a ten-task natural-language capability suite (T1--T10)
covering referent resolution, kind/state splitting, mid-flight
re-commanding, partial cancel, conditional trigger, route constraint,
formation, time-sequencing, and failure recovery, all passing; (ii) a
live LLM session of eight free-form commands executed end-to-end; (iii)
repeated DeepSeek-V4-Pro trials on three tactical scenarios (N=5 per
scenario) showing low and stable LLM turn / token cost per command. We
additionally report a cross-paradigm cost reference against OpenRA-RL, a
per-unit atomic baseline, to characterise the architectural cost of
atomic action APIs. Recent MCP-based UAV work shows that natural
language to MCP-mediated multi-entity execution is emerging in robotics;
this paper therefore claims the game/RTS branch of that design space
rather than MCP swarm control in general. We position
\texttt{openra\_mcp} as the first engine-FSM-minimal MCP co-pilot for a
production RTS --- distinct from per-unit atomic-API designs used by
both autonomous-player systems and an earlier open-source co-pilot
effort (HOPPINZQ, 2025) --- distinguish it from autonomous LLM-as-player
systems, robotics swarm-control systems, and broad-surface co-pilots in
turn-based or simulation games, and outline limitations (no formal user
study) and a future-work path toward small local low-latency models
distilled from LLM demonstrations.

\subsection{1. Introduction}\label{introduction}

The mechanical cost of issuing a tactical decision in a real-time
strategy (RTS) game often dwarfs the cost of making it. A player may
decide in one second to split a force, route the armour through the
middle, send the APCs around the flanks, delay the raiders by five
seconds, and cancel only the third group; expressing that decision
through selection boxes, control groups, queued orders, and rapid camera
movement takes far longer and is error-prone, especially in large
engagements with dozens of mixed units. This paper asks whether plain
natural language --- the same words the player would use to describe the
plan to a friend --- can serve as the tactical control interface for a
production RTS, and what kind of engine-side abstraction makes that
practical.

Most prior LLM-and-RTS research takes the \textbf{LLM-as-player} path:
substituting the human with an agent and evaluating on game outcome.
AlphaStar reached grandmaster StarCraft II play with reinforcement
learning (Vinyals et al., 2019); recent LLM-based agents extend this
line into language-driven strategic planning (Ma et al., 2024; Shao et
al., 2024; Wang et al., 2023). Even systems that accept a
natural-language prompt from a present human --- most prominently HIVE
(Anne et al., 2025) --- typically run plan-once: the human issues a
single instruction, the LLM emits a static plan, and the agents execute
it autonomously until the engagement ends.

We pursue a different question --- the \textbf{LLM-as-co-pilot} path ---
and a more specific design target: a \emph{plain-natural-language}
tactical interface for a real-time game with dozens of units, in which
the player stays in the loop and the LLM is restricted to translation.
The player never names actor IDs, computes coordinates, or selects from
a fifty-tool menu; they speak the way they would describe the maneuver
to a teammate. This direction has matured in code editing (GitHub
Copilot, 2021; Cursor, 2023; Claude Code) and in flight simulators
(Microsoft Flight Simulator Copilot, 2024), and exists in a pre-LLM form
in games --- Tom Clancy's EndWar (Ubisoft, 2008) demonstrated
voice-commanded RTS with rule-based speech recognition more than fifteen
years ago. At the same time, adjacent robotics work has begun to connect
LLMs, MCP gateways, and drone systems, including UAV swarm mission
execution (Ramos-Silva \& Burke, 2026; Iannoli et al., 2026). We
therefore do not claim to introduce MCP-mediated multi-entity control in
general. The gap we target is narrower: to our knowledge, plain-language
mixed-initiative co-piloting has not been instantiated for a production
RTS engine with a real AI opponent using a modern LLM, and existing
MCP-based game co-pilots target either turn-based games (STS2MCP for
Slay the Spire 2; Gennadiyev, 2025) or simulation/sandbox games
(oni\_mcp for Oxygen Not Included; LightJUNction, 2026) --- not
real-time strategy.

We present \texttt{openra\_mcp}, a prototype mixed-initiative co-pilot
for OpenRA, an open-source engine that recreates and modernizes classic
RTS games such as Red Alert and Command \& Conquer (OpenRA Project,
n.d.). The player issues natural-language commands such as ``send the
tanks down the middle and have the APCs flank from both sides'' or
``send the front group first, then launch the raiders five seconds
later.'' The LLM maps the command to structured Model Context Protocol
(MCP) tool calls. MCP provides a standard way to connect AI assistants
to external tools and data systems (Anthropic, 2024); in this prototype,
it is used as the bridge between the LLM client and an OpenRA unit
control interface. The OpenRA engine then handles pathfinding,
collision, attack-move behavior, and unit autonomy.

The current system is deliberately small on the engine side. After
ablation, the LLM-facing game interface exposes two squad execution
primitives: \texttt{Assault}, which pushes a unit set toward a target
cell or actor, and \texttt{Protection}, which holds or defends a
position. Higher-level tactics are not encoded as many specialized
engine-side finite-state machines. They are composed above the engine
boundary through Python/LLM logic and issued through MCP calls only when
the tactical state changes.

This paper makes the following contributions:

\begin{enumerate}
\def\labelenumi{\arabic{enumi}.}
\tightlist
\item
  \textbf{Plain natural language as the tactical control interface for a
  production RTS.} We demonstrate that a player can drive 30 mixed units
  through free-form natural-language commands --- without naming actor
  IDs, computing coordinates, or selecting from a large tool menu ---
  and have the engine execute the intended maneuver. To our knowledge
  this is the first such game/RTS co-pilot demonstration with an
  engine-FSM-minimal MCP surface on a real RTS engine with a real AI
  opponent --- distinct from per-unit atomic-API designs in
  autonomous-player systems (OpenRA-RL Contributors, 2026) and an
  earlier open-source co-pilot effort (HOPPINZQ, 2025).
\item
  \textbf{Two-primitive engine-FSM design.} An ablation-justified
  game-side interface --- \texttt{Assault} and \texttt{Protection} ---
  sufficient to compose ten distinct natural-language tactical
  capabilities (referent resolution, kind/state splitting, mid-flight
  re-commanding, partial cancel, conditional trigger, route constraint,
  formation, time-sequencing, failure recovery), with the higher-level
  tactics composed on the LLM side rather than encoded as specialized
  engine FSMs.
\item
  \textbf{Paradigm distinction within the LLM-game and MCP-control
  space.} We separate \emph{LLM-as-player} (autonomous game outcome)
  from \emph{LLM-as-co-pilot} (human-in-loop intent translation) and
  further distinguish two co-pilot sub-philosophies --- broad-surface
  (e.g.~oni\_mcp, \textasciitilde320 tools for Oxygen Not Included)
  versus engine-FSM-minimal (this work, 17 tools plus 2 FSM primitives)
  --- while explicitly separating game/RTS co-piloting from adjacent
  MCP-mediated robotics swarm control.
\item
  \textbf{Multi-axis evaluation on the same LLM.} A ten-task NL
  capability suite passing 10/10, a live LLM session of eight free-form
  commands, and repeated DeepSeek-V4-Pro trials on three tactical
  scenarios (N=5 per scenario) with per-run logs.
\item
  \textbf{Cross-paradigm cost reference.} Under the same LLM, the
  per-unit atomic baseline (OpenRA-RL) requires 40 ± 7 LLM responses and
  755 ± 218k tokens (N=3 full games) to reach a terminal state; this is
  provided as an architectural cost reference, not a head-to-head
  gameplay benchmark.
\item
  \textbf{Open and reproducible artifacts.} Engine fork, MCP server,
  capability suite, runners, and raw per-turn logs are released. To our
  knowledge this is the first open-source RTS co-pilot accompanied by a
  paper, ablation, and multi-axis repeated evaluation; an earlier
  open-source effort (HOPPINZQ, 2025) adopts a per-unit atomic API
  without a paper or controlled evaluation.
\item
  \textbf{A future-work path} toward small local models distilled from
  high-quality LLM demonstrations, emitting low-latency MCP command JSON
  to reduce per-command LLM latency for real-time play.
\end{enumerate}

\subsection{2. Related Work}\label{related-work}

Prior research on language-driven RTS systems falls into two distinct
paradigms that share surface mechanics but optimize for different
objectives:

\begin{itemize}
\tightlist
\item
  \textbf{Paradigm A --- LLM-as-player.} A language model (alone or
  hybridized with RL) replaces the human player and is evaluated on
  game-outcome metrics (win rate, score, reward). AlphaStar,
  TextStarCraft II, SwarmBrain, HIMA, OpenRA-RL, and HIVE all sit in
  this paradigm.
\item
  \textbf{Paradigm B --- LLM-as-co-pilot.} A language model translates a
  present human player's natural-language intent into engine-level
  actions; the human retains strategy, economy, and battlefield
  interpretation. Evaluation targets the translation cost (turns,
  tokens, latency) and intent fidelity, not game outcome.
\end{itemize}

\texttt{openra\_mcp} is positioned in Paradigm B. Conflating the two
paradigms invites mismatched evaluation criteria (e.g.~asking a co-pilot
for a win rate, or asking an autonomous agent for prompt-fidelity); we
treat all Paradigm A systems as comparison points on the \emph{cost}
axis only, not as precedents.

\subsubsection{2.0 Game + MCP landscape (May
2026)}\label{game-mcp-landscape-may-2026}

We surveyed open-source MCP servers connected to interactive games via
GitHub (keyword \texttt{mcp\ game}, sorted by stars, ≥40 stars retained,
25 repositories examined). The results cluster into four groups:

\begin{longtable}[]{@{}
  >{\raggedright\arraybackslash}p{(\linewidth - 6\tabcolsep) * \real{0.2308}}
  >{\raggedright\arraybackslash}p{(\linewidth - 6\tabcolsep) * \real{0.2308}}
  >{\raggedleft\arraybackslash}p{(\linewidth - 6\tabcolsep) * \real{0.3077}}
  >{\raggedright\arraybackslash}p{(\linewidth - 6\tabcolsep) * \real{0.2308}}@{}}
\toprule\noalign{}
\begin{minipage}[b]{\linewidth}\raggedright
Group
\end{minipage} & \begin{minipage}[b]{\linewidth}\raggedright
Representative repositories
\end{minipage} & \begin{minipage}[b]{\linewidth}\raggedleft
Stars (≥40 each)
\end{minipage} & \begin{minipage}[b]{\linewidth}\raggedright
Position vs \texttt{openra\_mcp}
\end{minipage} \\
\midrule\noalign{}
\endhead
\bottomrule\noalign{}
\endlastfoot
(i) Game-development tools for designers &
\texttt{Coding-Solo/godot-mcp} (3.8k), \texttt{ee0pdt/Godot-MCP} (570),
\texttt{AnkleBreaker-Studio/unity-mcp-server} (199, 268 tools),
\texttt{tugcantopaloglu/godot-mcp} (219, 149 tools),
\texttt{MubarakHAlketbi/game-asset-mcp} (137),
\texttt{youichi-uda/unity-mcp-pro-plugin} (57, 147 tools),
\texttt{HurtzDonutStudios/ai-forge-mcp} (61, 565 tools) & --- &
Different user (developer, not player); orthogonal \\
(ii) Autonomous in-game agents (Paradigm A) &
\texttt{Gennadiyev/STS2MCP} (363, Slay the Spire 2),
\texttt{CharTyr/STS2-Agent} (243),
\texttt{yuniko-software/minecraft-mcp-server} (586),
\texttt{notpoiu/roblox-executor-mcp} (60),
\texttt{Whale-io/lets-play-a-game} (132) & --- & Same MCP
infrastructure, different paradigm \\
(iii) Emulator / debug / reverse engineering & \texttt{drhelius/Gearboy}
(1.1k, Game Boy emulator + MCP debug), \texttt{drhelius/Gearsystem}
(369), \texttt{0xhackerfren/frida-game-hacking-mcp} (65),
\texttt{bethington/cheat-engine-server-python} (49) & --- & Non-gameplay
infra \\
(iv) Co-pilot for human players (Paradigm B) &
\texttt{LightJUNction/OniMods/oni\_mcp} (Oxygen Not Included),
\texttt{openra\_mcp} (this work) & 2 repositories total & Direct
paradigm peers \\
\end{longtable}

Group (iv) --- the design space this paper targets --- is sparse. The
two repositories represent two distinguishable sub-philosophies within
Paradigm B:

\begin{itemize}
\tightlist
\item
  \textbf{Broad-surface co-pilot.} \texttt{oni\_mcp} exposes
  approximately 320 tools spanning colony state, building configuration,
  scheduling, automation, research, and sandbox controls, with a small
  DSL-style batch wrapper (\texttt{agent\_program\_execute}). The LLM is
  expected to navigate a wide tool catalogue and decompose tactics
  itself. This fits Oxygen Not Included's low time pressure (a
  base-building simulation), where LLM latency does not block gameplay.
\item
  \textbf{Engine-FSM-minimal co-pilot.} \texttt{openra\_mcp} exposes 17
  tools, of which two (\texttt{spawn\_squad},
  \texttt{spawn\_squad\_batch}) trigger engine-side FSMs that execute
  per-tick movement, autotargeting, and cohesion internally. The LLM
  issues a task-level command, then the engine runs the maneuver without
  further LLM calls until the player or an event changes the intent.
  This fits real-time strategy, where per-tick LLM control would exceed
  the game's clock rate.
\end{itemize}

These philosophies are complementary rather than competing --- they suit
different game-time-pressure regimes. We adopt the engine-FSM-minimal
approach because the target game is real-time and the per-tick budget
for LLM round-trips is effectively zero.

\subsubsection{2.1 Cross-domain MCP and swarm-control
precedents}\label{cross-domain-mcp-and-swarm-control-precedents}

Recent robotics work provides an important boundary condition for this
paper: natural language to MCP-mediated actuation is no longer unique to
games. Ramos-Silva and Burke (2026) propose a universal LLM drone
command-and-control interface using MCP over MAVLink/MAVSDK for real and
simulated drones. Even closer, Iannoli et al.~(2026) present a
Web-of-Drones framework in which a user states a UAV swarm mission in
natural language and an LLM Agent Core interacts through an MCP gateway
and Web-of-Things abstractions to execute area coverage, formation, and
smart-irrigation missions in ArduPilot-based simulation.

These systems are strong adjacent precedents, and they rule out a broad
claim that \texttt{openra\_mcp} is the first LLM/MCP system for
controlling many entities. Our claim is narrower and game-specific.
\texttt{openra\_mcp} targets a production RTS rather than a robotics
swarm simulator; the controlled entities are heterogeneous combat units
rather than UAV platforms; the user is an active player who retains
strategy, economy, and battlefield judgment rather than a mission
designer who leaves the system to complete the task; and the control
layer emits squad-level game primitives rather than physical-device
actions. The relationship is therefore complementary: robotics work
establishes a cross-domain MCP multi-entity precedent, while this paper
studies the RTS game-control branch of the same larger design space.

\subsubsection{2.2 Cross-domain co-pilot precedents (Paradigm B
lineage)}\label{cross-domain-co-pilot-precedents-paradigm-b-lineage}

The mature reference point for LLM co-pilots is not games at all but
software engineering. GitHub Copilot, Cursor, and Claude Code translate
developer natural-language intent into code edits and tool calls while
the human retains goal selection, design judgment, and final acceptance
authority. The same pattern appears in Microsoft Flight Simulator's
voice copilot demonstration (2024) for cockpit operations and in
OpenAI's ChatGPT Code Interpreter for data-analysis workflows. These
systems share three properties with \texttt{openra\_mcp}: (i) the human
stays in the loop and owns the goal, (ii) the LLM emits structured tool
calls auditable by the human, and (iii) the underlying engine ---
compiler, OS, simulator --- owns the execution loop. To our knowledge,
this paradigm has not previously been instantiated against a production
RTS engine with a real opponent AI.

\subsubsection{2.3 NL game control before
LLMs}\label{nl-game-control-before-llms}

Pre-LLM RTS work demonstrated that natural-language unit control is
desirable but constrained by available NL technology. Tom Clancy's
EndWar (Ubisoft, 2008) used speech recognition with a fixed grammar to
issue squad-level commands (``Bravo, attack hostile two'') and remains
the most prominent commercial example of voice-controlled RTS. Black \&
White (Lionhead, 2001) used gesture and voice as primary input for a god
game. These systems established the design intent --- humans command,
the engine executes --- but were limited to rule-based parsers.
\texttt{openra\_mcp} continues this lineage with modern LLM-based
natural-language understanding.

\subsubsection{2.4 Paradigm A --- LLM-as-player in RTS (compared, not
precedent)}\label{paradigm-a-llm-as-player-in-rts-compared-not-precedent}

AlphaStar reached elite human performance in StarCraft II via
reinforcement learning, substituting the human player with an agent
(Vinyals et al., 2019). TextStarCraft II converts StarCraft II into a
text-based benchmark for evaluating LLM strategic decision-making and
introduces a chain-of-summarization approach for real-time strategic
settings (Ma et al., 2024). SwarmBrain uses an LLM-powered high-level
module with a lower-level reflex layer for StarCraft II, separating
strategic reasoning from faster tactical reactions (Shao et al., 2024).
HIMA extends this line through a hierarchical imitation multi-agent
framework for StarCraft II, with specialized imitation agents
orchestrated by a strategic planner (Ahn et al., 2025). OpenRA-RL frames
Red Alert as a platform for scripted bots, reinforcement learning, and
LLM players via a 48-tool atomic action API (OpenRA-RL Contributors,
2026); we use it as our direct cost-axis baseline in §6. An earlier
open-source RTS effort, \texttt{hoppinai-mcp-red95} (HOPPINZQ, 2025),
adopts the same per-unit atomic API style but in Paradigm B (co-pilot)
mode --- illustrating that the atomic-vs-task-level architectural choice
is orthogonal to the autonomy axis.

HIVE deserves separate treatment because it is frequently miscategorized
as a co-pilot system. HIVE enables a single human to coordinate swarms
of up to 2,000 agents through natural-language dialogue with an LLM,
generating behavior-tree-like plans for a custom multi-agent benchmark
(Anne et al., 2025). The architecture, however, is plan-once: the LLM is
invoked once per human command and emits a static plan executed for the
remainder of the engagement. The agents are JAX-vectorized abstract
particles rather than RTS units with pathfinding, weapon systems, or a
learning AI opponent; the enemy side is hardcoded. The code for the
LLM-to-DSL planner and the benchmark engine is not publicly released at
the time of writing. HIVE is therefore better characterized as an
\emph{adjacent} Paradigm A demonstration that uses NL prompts as a
configuration interface to a planner, rather than as a closed-loop
co-pilot for a production RTS. \texttt{openra\_mcp} differs along all of
these axes: production engine, closed-loop event-driven LLM invocation,
real AI opponent, open-source implementation.

\subsubsection{2.5 Engine-level autonomy
abstractions}\label{engine-level-autonomy-abstractions}

Behavior trees and hierarchical task networks are the canonical
engine-level abstractions for decomposing high-level commands into
per-tick behavior (Champandard, 2007). HIVE's five behavior trees follow
this lineage; our two-primitive squad FSM (Assault + Protection) is a
deliberately minimal specialization of the same idea, validated
empirically by ablation showing that ten natural-language tactical
capabilities (T1-T10) compose from these two primitives alone (§6.4).
Real-time human-AI coordination systems such as HLA address LLM latency
by separating slow language reasoning, faster macro-action generation,
and reactive execution (Liu et al., 2023). Our architecture aligns with
this principle: the LLM should not sit in a high-frequency loop; it
should translate intent and compose task-level actions while lower
layers handle per-tick execution.

\subsubsection{2.6 Mixed-initiative HCI and tool-use
frameworks}\label{mixed-initiative-hci-and-tool-use-frameworks}

The mixed-initiative framing comes from earlier HCI work arguing that
useful interfaces require careful sharing of initiative between humans
and agents rather than simple automation (Horvitz, 1999). More recent
guidelines for human-AI interaction emphasize making clear what the
system can do, how uncertainty and failure are handled, and when the
user remains in control (Amershi et al., 2019). \texttt{openra\_mcp}
applies these ideas to RTS control by assigning strategic judgment,
economy, and battlefield interpretation to the player while assigning
tactical translation to the LLM. Embodied LLM agents such as Voyager
demonstrate the power of tool-using language models in game-like
environments (Wang et al., 2023). The Model Context Protocol (Anthropic,
2024) provides the tool-call infrastructure used by \texttt{openra\_mcp}
to expose the squad primitives to the LLM. The resulting research
question is therefore not ``Can an LLM play Red Alert?'' but ``Can
natural language become a reliable tactical control layer for a human
player of a real-time strategy game?''

\subsubsection{2.7 Position summary}\label{position-summary}

\begin{longtable}[]{@{}
  >{\raggedright\arraybackslash}p{(\linewidth - 10\tabcolsep) * \real{0.1667}}
  >{\raggedright\arraybackslash}p{(\linewidth - 10\tabcolsep) * \real{0.1667}}
  >{\raggedright\arraybackslash}p{(\linewidth - 10\tabcolsep) * \real{0.1667}}
  >{\raggedright\arraybackslash}p{(\linewidth - 10\tabcolsep) * \real{0.1667}}
  >{\raggedright\arraybackslash}p{(\linewidth - 10\tabcolsep) * \real{0.1667}}
  >{\raggedright\arraybackslash}p{(\linewidth - 10\tabcolsep) * \real{0.1667}}@{}}
\toprule\noalign{}
\begin{minipage}[b]{\linewidth}\raggedright
Line of work
\end{minipage} & \begin{minipage}[b]{\linewidth}\raggedright
Paradigm
\end{minipage} & \begin{minipage}[b]{\linewidth}\raggedright
Environment
\end{minipage} & \begin{minipage}[b]{\linewidth}\raggedright
Human role
\end{minipage} & \begin{minipage}[b]{\linewidth}\raggedright
LLM control layer
\end{minipage} & \begin{minipage}[b]{\linewidth}\raggedright
Status as precedent for \texttt{openra\_mcp}
\end{minipage} \\
\midrule\noalign{}
\endhead
\bottomrule\noalign{}
\endlastfoot
AlphaStar (Vinyals et al., 2019) & A & StarCraft II & Spectator/opponent
& Learned full-game RL policy & Comparison only --- different
paradigm \\
TextStarCraft II / HIMA (Ma 2024; Ahn 2025) & A & Textual or structured
SC2 & Outside the loop & Strategic LLM planning & Comparison only ---
different paradigm \\
SwarmBrain (Shao et al., 2024) & A & StarCraft II & Outside the loop &
LLM strategy + reflex layer & Comparison only --- different paradigm \\
OpenRA-RL (Contributors, 2026) & A & OpenRA, atomic 48-tool API &
Outside the loop & Per-unit atomic actions & Direct \emph{cost-axis}
baseline in §6, not paradigm precedent \\
HIVE (Anne et al., 2025) & A (adjacent) & JAX swarm sim, plan-once &
Single NL prompt, then absent & LLM emits static DSL plan & Closest
surface, but plan-once + toy engine + closed source --- not a Paradigm B
precedent \\
Universal LLM Drone C2 (Ramos-Silva \& Burke, 2026) & Adjacent robotics
& Real/simulated drones & Operator gives drone command & MCP over
MAVLink/MAVSDK & MCP physical-control precedent, not game/RTS \\
Web-of-Drones swarm (Iannoli et al., 2026) & Adjacent robotics &
ArduPilot UAV swarm simulation & Gives mission objective, then absent
during execution & LLM Agent Core + MCP gateway + WoT & Closest non-game
MCP multi-entity analogue; not RTS/mixed-initiative game control \\
Copilot / Cursor / Claude Code & B (other domain) & Code editor & Owns
design, accepts edits & NL-to-edit + tool calls & Cross-domain
precedent \\
EndWar (Ubisoft, 2008) & B (pre-LLM) & Commercial RTS & Owns strategy,
voice commander & Rule-based ASR + grammar & Lineage precedent \\
Hoppinai (HOPPINZQ, 2025) & B & OpenRA & Owns strategy, issues per-unit
commands & Per-unit atomic API (move/attack by actor\_id) & Earlier
open-source RTS co-pilot effort; atomic-API design choice
(cf.~OpenRA-RL); no paper or controlled evaluation \\
\textbf{\texttt{openra\_mcp} (this work)} & \textbf{B} & \textbf{OpenRA,
production RTS} & \textbf{Owns strategy/economy/judgment} &
\textbf{Two-primitive squad FSM + LLM composition} & \textbf{First
engine-FSM-minimal RTS co-pilot with paper, ablation, and multi-axis
evaluation} \\
\end{longtable}

\subsection{3. System Overview}\label{system-overview}

The system has four layers:

\begin{Shaded}
\begin{Highlighting}[]
\NormalTok{Human player intent}
\NormalTok{        {-}\textgreater{}}
\NormalTok{LLM tactical translator}
\NormalTok{        {-}\textgreater{}}
\NormalTok{MCP tool calls}
\NormalTok{        {-}\textgreater{}}
\NormalTok{OpenRA squad execution primitives}
\NormalTok{        {-}\textgreater{}}
\NormalTok{Engine unit autonomy}
\end{Highlighting}
\end{Shaded}

The human layer supplies strategic and tactical intent in natural
language. The LLM layer translates that intent into structured tool
calls. The MCP layer provides an auditable interface between the LLM
client and the game-side server. The OpenRA layer executes squad-level
primitives and delegates local movement and combat details to the
engine.

The design follows an explicit authority split:

\begin{itemize}
\tightlist
\item
  The player owns strategic judgment.
\item
  The player owns economy and production decisions.
\item
  The player owns battlefield interpretation.
\item
  The LLM owns tactical translation and composition.
\item
  The engine owns pathfinding, collision, and local unit behavior.
\end{itemize}

This authority split is central to the paper. The goal is not to make an
AI that plays the game for the user. The goal is to make an AI co-pilot
that turns the user's declared tactical intent into complex executable
control.

The architecture also creates an audit trail. A natural-language command
is not converted into invisible mouse input. It becomes a structured MCP
tool call, such as \texttt{spawn\_squad} or
\texttt{spawn\_squad\_batch}, with explicit unit identifiers, target
positions, and squad types. This matters for research reproducibility:
when a command succeeds or fails, the intermediate representation can be
inspected, compared across runs, or used later as a supervised training
target for smaller models.

\subsubsection{3.1 Example Translation}\label{example-translation}

The following simplified example illustrates the intended interface. A
player utterance is first interpreted at the tactical level and then
emitted as a structured MCP call:

\begin{Shaded}
\begin{Highlighting}[]
\FunctionTok{\{}
  \DataTypeTok{"player\_command"}\FunctionTok{:} \StringTok{"Tanks through the middle, APCs flank from both sides"}\FunctionTok{,}
  \DataTypeTok{"mcp\_tool"}\FunctionTok{:} \StringTok{"spawn\_squad\_batch"}\FunctionTok{,}
  \DataTypeTok{"arguments"}\FunctionTok{:} \FunctionTok{\{}
    \DataTypeTok{"squads"}\FunctionTok{:} \OtherTok{[}
      \FunctionTok{\{}
        \DataTypeTok{"squad\_type"}\FunctionTok{:} \StringTok{"Assault"}\FunctionTok{,}
        \DataTypeTok{"unit\_ids"}\FunctionTok{:} \OtherTok{[}\DecValTok{101}\OtherTok{,} \DecValTok{102}\OtherTok{,} \DecValTok{103}\OtherTok{]}\FunctionTok{,}
        \DataTypeTok{"target\_pos"}\FunctionTok{:} \FunctionTok{\{} \DataTypeTok{"x"}\FunctionTok{:} \DecValTok{60}\FunctionTok{,} \DataTypeTok{"y"}\FunctionTok{:} \DecValTok{50} \FunctionTok{\}}
      \FunctionTok{\}}\OtherTok{,}
      \FunctionTok{\{}
        \DataTypeTok{"squad\_type"}\FunctionTok{:} \StringTok{"Assault"}\FunctionTok{,}
        \DataTypeTok{"unit\_ids"}\FunctionTok{:} \OtherTok{[}\DecValTok{201}\OtherTok{,} \DecValTok{202}\OtherTok{,} \DecValTok{203}\OtherTok{]}\FunctionTok{,}
        \DataTypeTok{"target\_pos"}\FunctionTok{:} \FunctionTok{\{} \DataTypeTok{"x"}\FunctionTok{:} \DecValTok{60}\FunctionTok{,} \DataTypeTok{"y"}\FunctionTok{:} \DecValTok{40} \FunctionTok{\}}
      \FunctionTok{\}}\OtherTok{,}
      \FunctionTok{\{}
        \DataTypeTok{"squad\_type"}\FunctionTok{:} \StringTok{"Assault"}\FunctionTok{,}
        \DataTypeTok{"unit\_ids"}\FunctionTok{:} \OtherTok{[}\DecValTok{301}\OtherTok{,} \DecValTok{302}\OtherTok{,} \DecValTok{303}\OtherTok{]}\FunctionTok{,}
        \DataTypeTok{"target\_pos"}\FunctionTok{:} \FunctionTok{\{} \DataTypeTok{"x"}\FunctionTok{:} \DecValTok{60}\FunctionTok{,} \DataTypeTok{"y"}\FunctionTok{:} \DecValTok{60} \FunctionTok{\}}
      \FunctionTok{\}}
    \OtherTok{]}
  \FunctionTok{\}}
\FunctionTok{\}}
\end{Highlighting}
\end{Shaded}

The exact unit identifiers and target cells come from the current game
state and the composition layer. OpenRA then executes the squad orders
through its own pathfinding, collision, and combat systems. The LLM does
not issue per-tick movement commands.

\subsection{4. Two Execution Primitives}\label{two-execution-primitives}

Early versions of the project explored a larger daemon and DSL surface.
After the v2 tactical capability suite, the implementation was reduced.
The current LLM-facing design relies on two squad execution primitives.

\texttt{Assault} pushes a set of units toward a target cell or actor. It
is the core primitive for attacks, movement, pincer arms,
route-constrained pushes, and timed raids. \texttt{Protection} holds a
position and defends a target cell. These primitives are intentionally
execution-level, not strategic. They do not decide what the player
should build, whether a fight is favorable, or which strategic objective
matters.

Higher-level tactics are composed outside the engine primitive layer.
For example, pincer movement is represented as two Assault squads
approaching from different sides. A time-sequenced attack is represented
as one squad launched first and a second squad launched after a delay. A
feint-and-raid plan is represented as separate squad calls with
different targets and timing.

This design keeps the game-side interface small and makes the LLM's
output auditable: each player command can be inspected as a sequence of
MCP calls over a small vocabulary.

\subsection{5. Tactical Composition From Natural
Language}\label{tactical-composition-from-natural-language}

The prototype demonstrates several forms of tactical composition:

\begin{itemize}
\tightlist
\item
  referential control: ``the left group'' versus ``the right group'';
\item
  kind-based splitting: tanks through the middle, APCs on the flanks;
\item
  state-based splitting: damaged units return, healthy units continue;
\item
  mid-flight recommanding: stop one ongoing movement and replace it;
\item
  partial cancellation: recall one subgroup while others continue;
\item
  conditional behavior: retreat if a condition is met;
\item
  route constraints: avoid the center and move through the side;
\item
  formation-like behavior: tanks ahead, APCs behind;
\item
  time sequencing: launch the main force first, raiders later;
\item
  failure recovery: re-plan when a group is stuck.
\end{itemize}

These are not merely textual suggestions. The LLM produces tool calls
that the OpenRA bridge executes in the game engine.

This section is also the bridge to the local-model future work. The
target output is not an unstructured essay and not a hidden policy. It
is a compact JSON-like action specification over a small tool
vocabulary. That makes the problem suitable for supervised fine-tuning:
large models can generate high-quality demonstrations, while smaller
local models can be trained to imitate the mapping from a compact
game-state summary plus player instruction to MCP calls.

\subsection{6. Evaluation}\label{evaluation}

The current evidence base contains four parts: a natural-language
tactical capability suite, a live LLM demonstration, earlier
two-primitive tactical demonstrations, and an ablation showing that
unused daemon/DSL machinery could be removed without losing the
demonstrated tactical path.

\subsubsection{6.1 Natural-Language Tactical Capability
Suite}\label{natural-language-tactical-capability-suite}

The v2 capability suite tests ten tactical capabilities that are
difficult to express as a single conventional RTS command. Each scenario
begins with an English or Chinese natural-language instruction and
evaluates whether the system selects the correct units, decomposes the
task, and satisfies the tactical intent in OpenRA. Table 1 reports the
final merged result across the base post-ablation run and targeted retry
files. This is intentionally reported as a merged final result: the raw
\texttt{v2\_post\_ablation.csv} file alone contains threshold/scenario
failures for T3, T5, T7, and T8, and the corrected outcomes come from
\texttt{v2\_retry\_4fails.csv} and \texttt{v2\_t7\_retry2.csv}.

\begin{longtable}[]{@{}
  >{\raggedright\arraybackslash}p{(\linewidth - 12\tabcolsep) * \real{0.1250}}
  >{\raggedright\arraybackslash}p{(\linewidth - 12\tabcolsep) * \real{0.1250}}
  >{\raggedright\arraybackslash}p{(\linewidth - 12\tabcolsep) * \real{0.1250}}
  >{\raggedleft\arraybackslash}p{(\linewidth - 12\tabcolsep) * \real{0.1667}}
  >{\raggedleft\arraybackslash}p{(\linewidth - 12\tabcolsep) * \real{0.1667}}
  >{\raggedright\arraybackslash}p{(\linewidth - 12\tabcolsep) * \real{0.1250}}
  >{\raggedleft\arraybackslash}p{(\linewidth - 12\tabcolsep) * \real{0.1667}}@{}}
\toprule\noalign{}
\begin{minipage}[b]{\linewidth}\raggedright
ID
\end{minipage} & \begin{minipage}[b]{\linewidth}\raggedright
Capability
\end{minipage} & \begin{minipage}[b]{\linewidth}\raggedright
Natural-language command
\end{minipage} & \begin{minipage}[b]{\linewidth}\raggedleft
Units
\end{minipage} & \begin{minipage}[b]{\linewidth}\raggedleft
Subtasks
\end{minipage} & \begin{minipage}[b]{\linewidth}\raggedright
Final result
\end{minipage} & \begin{minipage}[b]{\linewidth}\raggedleft
Latency
\end{minipage} \\
\midrule\noalign{}
\endhead
\bottomrule\noalign{}
\endlastfoot
T1 & Referential resolution & ``Let the left group flank; the right
group stays still'' & 80 & 1 & Pass & 348 ms \\
T2 & Kind-based split & ``Tanks through the middle, APCs flank from both
sides'' & 100 & 3 & Pass & 356 ms \\
T3 & State-based split & ``Damaged units return; the rest keep pushing''
& 64 & 2 & Pass & 360 ms \\
T4 & Mid-flight re-command & ``Everyone push, stop, then switch to
two-side pincer'' & 100 & 3 & Pass & 8,736 ms \\
T5 & Partial cancel & ``The third group returns; the others continue'' &
64 & 6 & Pass & 3,636 ms \\
T6 & Conditional retreat & ``If enemy main force appears, retreat to the
bridge; otherwise continue'' & 75 & 2 & Pass & 1,712 ms \\
T7 & Route constraint & ``Group A goes directly; group B goes far right
first'' & 64 & 2 & Pass & 25,175 ms \\
T8 & Formation-like movement & ``Tanks in front, APCs behind; do not
clump'' & 64 & 2 & Pass & 310 ms \\
T9 & Time sequencing & ``Main force first; raiders attack from the side
five seconds later'' & 75 & 2 & Pass & 5,362 ms \\
T10 & Failure recovery & ``If a group gets stuck, re-plan the route'' &
75 & 1 & Pass & 20,657 ms \\
\end{longtable}

The final merged suite reaches 10/10 pass. The longer latencies in T4,
T7, and T10 reflect multi-stage or waiting behavior rather than only the
initial tool call time.

\subsubsection{6.2 Live LLM Demonstration}\label{live-llm-demonstration}

The live demonstration records Claude receiving player natural language,
reading game state, selecting units, and emitting \texttt{spawn\_squad}
or \texttt{spawn\_squad\_batch} calls. The demonstration contains eight
consecutive commands with no failed command in the recorded sequence.

\begin{longtable}[]{@{}
  >{\raggedleft\arraybackslash}p{(\linewidth - 4\tabcolsep) * \real{0.4000}}
  >{\raggedright\arraybackslash}p{(\linewidth - 4\tabcolsep) * \real{0.3000}}
  >{\raggedright\arraybackslash}p{(\linewidth - 4\tabcolsep) * \real{0.3000}}@{}}
\toprule\noalign{}
\begin{minipage}[b]{\linewidth}\raggedleft
Step
\end{minipage} & \begin{minipage}[b]{\linewidth}\raggedright
Player intent
\end{minipage} & \begin{minipage}[b]{\linewidth}\raggedright
Capability shown
\end{minipage} \\
\midrule\noalign{}
\endhead
\bottomrule\noalign{}
\endlastfoot
1 & Move all units to the lower-right area & 100-unit single-squad
movement \\
2 & Send APCs to upper-left while one tank stays still & Unit-kind split
and local exception \\
3 & Send the left half downward and the right half upward & Spatial
reference resolution \\
4 & Regroup at center, then split into four teams to four corners &
Multi-stage regroup and four-way split \\
5 & Cycle four teams around the map for 180 seconds & Long-running
composed patrol \\
6 & Regroup all units at lower-left & Global recall/regroup \\
7 & Send a 20-unit team first, then 80 units after eight seconds &
Time-sequenced small/large force coordination \\
8 & Execute a pincer around the central building & 50/50 split and
converging attack arms \\
\end{longtable}

This sequence is important because it uses a live LLM rather than a
fixed script: the model reads \texttt{get\_state}, selects unit IDs,
computes or chooses targets, and emits batch JSON.

\subsubsection{6.3 Two-Primitive
Demonstrations}\label{two-primitive-demonstrations}

Earlier E7 demonstrations validate the engine-side primitive path. In
those runs, \texttt{spawn\_squad\_batch} dispatched four squads in 31-34
ms, an eight-squad batch in about 10 ms, and a single 60-unit squad in
about 18 ms. Demonstrated tactics include single-squad assault,
four-corner batch movement, composed patrol, pincer movement, feint plus
raid, mixed APC/tank squads, combined-arms offset movement, large-unit
movement, and eight-direction batch movement.

These demonstrations support the architectural claim that complex
tactics do not require a large set of engine-side tactical FSMs. A small
execution layer can support richer behavior when the composition logic
lives above it.

\subsubsection{6.4 Ablation and Telemetry}\label{ablation-and-telemetry}

The phase ablation removed the unused daemon and DSL path while
retaining the \texttt{spawn\_squad} / \texttt{spawn\_squad\_batch}
execution route. The active MCP surface was reduced from 31 tools to 17
tools, and the Python implementation shrank by approximately 6,700
lines. The post-ablation capability path remained viable, which supports
the ``two engine primitives plus LLM-side composition'' design.

Development-session telemetry provides a rough observability check
rather than a controlled user-study result. Across existing
\texttt{decisions.jsonl} logs, \texttt{paper\_metrics.py} found 123
decision records, 117 successful decisions, 6 errors, 98 commands with
recorded natural-language input, and 92 unique natural-language inputs.
Across Claude transcript logs for this project, it found 781 OpenRA MCP
tool events. The median game-side decision latency was 140 ms, the
median logged LLM latency field was 600 ms, and the median time from a
user text turn to the first OpenRA MCP tool event was about 7.19
seconds. Because these logs include development and debugging sessions,
these numbers should be read as prototype telemetry, not as final
benchmark results.

\subsubsection{6.5 Repeated DeepSeek-V4-Pro tactical trials (N=5 per
scenario)}\label{repeated-deepseek-v4-pro-tactical-trials-n5-per-scenario}

To probe how stable the LLM translation cost is across repetitions and
how it scales with roster size, we re-ran three tactical scenarios five
times each at two roster sizes --- 30 units (12 e1 + 8 3tnk + 6 v2rl + 4
apc) and 100 units (40 e1 + 30 3tnk + 20 v2rl + 10 apc) --- through a
same-game OpenRA session driven by DeepSeek-V4-Pro: \emph{scen1} ---
push all mobile units to map bottom-right (78, 85); \emph{scen2} ---
split into four squads, push to four corners; \emph{scen3} --- two-phase
50/50 pincer onto map centre. Between runs, surviving units were
recalled to a rally point near the player base via a one-shot
\texttt{spawn\_squad}. Each scenario carries its own verification window
before unit positions are checked; the verify radius and per-corner unit
threshold scale with √(roster / 30) to account for on-arrival congestion
at larger roster sizes.

\textbf{30-unit trials (N=5 per scenario):}

\begin{longtable}[]{@{}
  >{\raggedright\arraybackslash}p{(\linewidth - 8\tabcolsep) * \real{0.1875}}
  >{\raggedleft\arraybackslash}p{(\linewidth - 8\tabcolsep) * \real{0.2500}}
  >{\raggedright\arraybackslash}p{(\linewidth - 8\tabcolsep) * \real{0.1875}}
  >{\raggedright\arraybackslash}p{(\linewidth - 8\tabcolsep) * \real{0.1875}}
  >{\raggedright\arraybackslash}p{(\linewidth - 8\tabcolsep) * \real{0.1875}}@{}}
\toprule\noalign{}
\begin{minipage}[b]{\linewidth}\raggedright
Scenario
\end{minipage} & \begin{minipage}[b]{\linewidth}\raggedleft
Success
\end{minipage} & \begin{minipage}[b]{\linewidth}\raggedright
LLM turns mean ± std (min, max)
\end{minipage} & \begin{minipage}[b]{\linewidth}\raggedright
Total tokens mean ± std (min, max)
\end{minipage} & \begin{minipage}[b]{\linewidth}\raggedright
Wallclock mean ± std (s)
\end{minipage} \\
\midrule\noalign{}
\endhead
\bottomrule\noalign{}
\endlastfoot
scen1 full push BR & 4/5 & 3.0 ± 0.0 (3, 3) & 7,136 ± 106 (6,999, 7,296)
& 77.9 ± 4.8 \\
scen2 4-corner split & 4/5 & 2.6 ± 0.89 (1, 3) & 7,009 ± 3,360 (1,028,
8,958) & 108.4 ± 54.9 \\
scen3 50/50 pincer & 2/5 & 5.6 ± 2.51 (2, 8) & 20,502 ± 11,535 (5,182,
35,038) & 206.5 ± 102.7 \\
\end{longtable}

\textbf{100-unit trials (N=5 per scenario, 3.3× larger roster):}

\begin{longtable}[]{@{}
  >{\raggedright\arraybackslash}p{(\linewidth - 8\tabcolsep) * \real{0.1875}}
  >{\raggedleft\arraybackslash}p{(\linewidth - 8\tabcolsep) * \real{0.2500}}
  >{\raggedright\arraybackslash}p{(\linewidth - 8\tabcolsep) * \real{0.1875}}
  >{\raggedright\arraybackslash}p{(\linewidth - 8\tabcolsep) * \real{0.1875}}
  >{\raggedright\arraybackslash}p{(\linewidth - 8\tabcolsep) * \real{0.1875}}@{}}
\toprule\noalign{}
\begin{minipage}[b]{\linewidth}\raggedright
Scenario
\end{minipage} & \begin{minipage}[b]{\linewidth}\raggedleft
Success
\end{minipage} & \begin{minipage}[b]{\linewidth}\raggedright
LLM turns mean ± std (min, max)
\end{minipage} & \begin{minipage}[b]{\linewidth}\raggedright
Total tokens mean ± std (min, max)
\end{minipage} & \begin{minipage}[b]{\linewidth}\raggedright
Wallclock mean ± std (s)
\end{minipage} \\
\midrule\noalign{}
\endhead
\bottomrule\noalign{}
\endlastfoot
scen1 full push BR & 0/5 & 3.8 ± 0.84 (3, 5) & 19,752 ± 10,195 (10,898,
35,881) & 154.4 ± 8.7 \\
scen2 4-corner split & 0/5 & 3.6 ± 0.55 (3, 4) & 18,407 ± 5,846 (11,776,
23,376) & 186.9 ± 3.4 \\
scen3 50/50 pincer & 0/5 & 5.8 ± 0.84 (5, 7) & 33,534 ± 11,836 (22,264,
50,086) & 305.0 ± 9.1 \\
\end{longtable}

\textbf{LLM-side cost scales sub-linearly with roster.} Roster grew
3.3×; LLM turns grew only 1.04--1.38× across the three scenarios. Token
usage grew 1.6--2.8×, the increase driven almost entirely by larger
observation payloads (longer unit lists in \texttt{get\_state}) rather
than by additional LLM round-trips.

\textbf{100-unit verification did not pass the success criterion, but
NL-to-action translation remained correct.} In every 100-unit trial, the
LLM produced syntactically valid tool calls matching the requested
tactic (e.g.~scen2 r01--r05 all emitted a four-squad batch with balanced
allocations of approximately 25 units each, and the engine reported
post-execution distributions of {[}14, 14, 26, 14{]} units per corner
--- 80 out of 100 units traversed correctly, but the per-corner ≥15-unit
threshold required all four corners filled). scen1 100-unit consistently
delivered 56/100 units to within radius 18 of the target, falling 4
units short of the 60\% threshold. scen3 100-unit suffered from
engine-level congestion during the phase-1-to-phase-2 transition,
delivering 0--28 units to the centre depending on traversal timing.
These are engine-side execution outcomes, not translation failures, and
we report them transparently rather than relaxing the criterion.

\textbf{30-unit per-scenario failure modes (N=15 total):} one DeepSeek
transient connection drop (scen2 r03), one DeepSeek malformed JSON in a
tool call (scen3 r01), one timing edge where the verify window fired
before units reached the target (scen1 r02), and two squad-overlap
interactions in scen3 (r02 / r03) where the LLM issued the phase-2
\texttt{spawn\_squad\_batch} without first calling
\texttt{cancel\_squad} on the phase-1 squad, leaving the phase-1 squad
still holding the units. The squad-overlap interaction is a real system
limitation documented in the project's design notes; in scen3 r04 / r05
the LLM issued explicit \texttt{cancel\_squad} calls between phases and
succeeded.

\textbf{Per-call patterns are stable across both roster sizes.} Each
scenario typically consists of one \texttt{get\_state} read, one
\texttt{spawn\_squad} (or \texttt{spawn\_squad\_batch}) dispatch, and an
optional \texttt{wait} or \texttt{done}. scen3 adds one extra dispatch
and one \texttt{wait} for the second phase. The LLM never fabricated a
tool name, never produced a tool argument the engine rejected as
malformed, and consistently produced balanced unit allocations across
split-target requests.

Raw per-run data:
\texttt{logs/rl\_compare/our\_deepseek\_results\_runs\_n30.csv} and
\texttt{logs/rl\_compare/our\_deepseek\_results\_runs\_n100.csv}.
Summaries: \texttt{our\_deepseek\_summary\_n30.csv},
\texttt{our\_deepseek\_summary\_n100.csv}.

\subsubsection{6.6 Cross-paradigm cost reference: per-unit atomic
baseline}\label{cross-paradigm-cost-reference-per-unit-atomic-baseline}

To characterise the architectural cost of atomic per-unit MCP APIs ---
not to make a head-to-head gameplay claim --- we ran DeepSeek-V4-Pro
through the OpenRA-RL agent harness (Contributors, 2026), which exposes
48 atomic tools (\texttt{move\_units}, \texttt{attack\_target},
\texttt{build\_unit}, \texttt{advance}, etc.) over a docker-hosted
OpenRA fork. Three independent full games were run.

\begin{longtable}[]{@{}
  >{\raggedright\arraybackslash}p{(\linewidth - 14\tabcolsep) * \real{0.1000}}
  >{\raggedleft\arraybackslash}p{(\linewidth - 14\tabcolsep) * \real{0.1333}}
  >{\raggedleft\arraybackslash}p{(\linewidth - 14\tabcolsep) * \real{0.1333}}
  >{\raggedleft\arraybackslash}p{(\linewidth - 14\tabcolsep) * \real{0.1333}}
  >{\raggedleft\arraybackslash}p{(\linewidth - 14\tabcolsep) * \real{0.1333}}
  >{\raggedleft\arraybackslash}p{(\linewidth - 14\tabcolsep) * \real{0.1333}}
  >{\raggedleft\arraybackslash}p{(\linewidth - 14\tabcolsep) * \real{0.1333}}
  >{\raggedright\arraybackslash}p{(\linewidth - 14\tabcolsep) * \real{0.1000}}@{}}
\toprule\noalign{}
\begin{minipage}[b]{\linewidth}\raggedright
Run
\end{minipage} & \begin{minipage}[b]{\linewidth}\raggedleft
LLM responses
\end{minipage} & \begin{minipage}[b]{\linewidth}\raggedleft
Tool calls
\end{minipage} & \begin{minipage}[b]{\linewidth}\raggedleft
Prompt tokens
\end{minipage} & \begin{minipage}[b]{\linewidth}\raggedleft
Completion tokens
\end{minipage} & \begin{minipage}[b]{\linewidth}\raggedleft
Total tokens
\end{minipage} & \begin{minipage}[b]{\linewidth}\raggedleft
Wallclock (s)
\end{minipage} & \begin{minipage}[b]{\linewidth}\raggedright
Outcome
\end{minipage} \\
\midrule\noalign{}
\endhead
\bottomrule\noalign{}
\endlastfoot
1 & 47 & 78 & 980,323 & 13,940 & 994,263 & 518.9 & LOSE \\
2 & 33 & 55 & 550,221 & 9,028 & 559,249 & 440.8 & LOSE \\
3 & 40 & 66 & 699,546 & 11,585 & 711,131 & 556.0 & LOSE \\
\textbf{mean ± std} & \textbf{40.0 ± 7.0} & \textbf{66.3 ± 11.5} &
\textbf{743,363 ± 218,540} & \textbf{11,518 ± 2,462} & \textbf{754,881 ±
217,924} & \textbf{505.2 ± 58.9} & \textbf{3/3 LOSE} \\
\end{longtable}

Average prompt size per response across runs is \textasciitilde18--21k
tokens, dominated by the accumulating conversation history and the
48-tool schema repeated on every call. Average latency per LLM call is
11.0--14.7 s. All three games reached LOSE without exhausting the
100-turn / 1800-second cap.

\begin{quote}
\textbf{Caveat --- this is a cost reference, not a benchmark.} The
OpenRA-RL run plays a full game (economy + combat + outcome), while our
\texttt{openra\_mcp} scenarios in §6.5 are isolated tactical commands on
a pre-trained roster. Player-owned economy is by design in
\texttt{openra\_mcp} (see §3), so the workloads are not
apples-to-apples; we therefore do not claim a direct gameplay or
efficiency win over OpenRA-RL. The numbers serve a different purpose:
they illustrate the per-turn prompt-size and total- token implications
of exposing per-unit atomic actions as the LLM's control surface, which
is the architectural choice this paper argues against for real-time
strategy.
\end{quote}

For context, a same-engine scripted baseline (no LLM, the official
\texttt{ScriptedBot} driving identical isolated tactical tasks) requires
3,369 ± 9 environment round-trips per scenario set across three runs.
This is an action-granularity baseline that quantifies how much repeated
control work an atomic engine API would push onto any controller --- LLM
or otherwise --- that has not been promoted to a task-level abstraction.

Raw logs and summaries: \texttt{logs/rl\_compare/}.

\subsubsection{6.7 Model robustness: smaller-LLM
ablation}\label{model-robustness-smaller-llm-ablation}

A natural concern with co-pilot systems built on large frontier models
is whether the underlying paradigm collapses when the model is
downgraded. To probe this, we re-ran the §6.5 N=5 protocol on the
30-unit roster with \texttt{deepseek-v4-flash}, DeepSeek's smaller /
cheaper thinking-mode variant, holding the system prompt, tool surface,
scenarios, verify windows, and roster identical to the V4-Pro baseline.

\begin{longtable}[]{@{}lll@{}}
\toprule\noalign{}
Metric & V4-Pro (baseline) & V4-Flash \\
\midrule\noalign{}
\endhead
\bottomrule\noalign{}
\endlastfoot
scen1 success & 4/5 & \textbf{5/5} \\
scen2 success & 4/5 & 4/5 \\
scen3 success & 2/5 & 0/5 \\
\textbf{Overall} & \textbf{10/15 (67\%)} & \textbf{9/15 (60\%)} \\
scen1 total tokens (mean ± std) & 7,136 ± 106 & 8,835 ± 2,897 \\
scen2 total tokens (mean ± std) & 7,009 ± 3,360 & 7,908 ± 1,603 \\
scen3 total tokens (mean ± std) & 20,502 ± 11,535 & 20,080 ± 8,479 \\
scen1 wallclock (s, mean ± std) & 77.9 ± 4.8 & 131.6 ± 3.0 \\
scen2 wallclock (s, mean ± std) & 108.4 ± 54.9 & 134.9 ± 66.4 \\
scen3 wallclock (s, mean ± std) & 206.5 ± 102.7 & 228.0 ± 119.8 \\
\end{longtable}

Three observations from this ablation:

\begin{enumerate}
\def\labelenumi{\arabic{enumi}.}
\tightlist
\item
  \textbf{Single-step tactics survive a model downgrade}. On scen1 (one
  \texttt{spawn\_squad}) and scen2 (one \texttt{spawn\_squad\_batch}),
  V4-Flash matches or exceeds V4-Pro. This supports the paradigm-level
  claim of the paper: when the engine-side FSM does the per-tick work,
  the LLM-side call surface is small enough that even a smaller
  thinking-mode model can drive it.
\item
  \textbf{Multi-stage tactics expose model capability gaps}. On scen3
  (phase-1 split, then phase-2 convergence), V4-Flash drops from 2/5 to
  0/5. The failure mode is uniform: V4-Flash issues the phase-2
  \texttt{spawn\_squad\_batch} without first calling
  \texttt{cancel\_squad} on the phase-1 squads, so the old squad
  assignments continue to hold the units. V4-Pro hits the same trap
  occasionally; V4-Flash hits it every time. This is a model-level
  reasoning gap, not a paradigm failure: the tool surface offers
  \texttt{cancel\_squad}, but the smaller model does not select it
  without more explicit prompting.
\item
  \textbf{Token cost is not the savings axis}. V4-Flash uses roughly the
  same number of tokens per scenario as V4-Pro because the prompt
  schema, tool definitions, and conversation history are identical. The
  savings from a smaller model come from per-token price and per-call
  latency, not from a shorter dialog.
\end{enumerate}

This ablation is consistent with the broader thesis that the paradigm
--- plain natural language plus an engine-side FSM --- is robust to
model choice for the common one-shot tactical step, and that the
residual fragility lives in multi-stage tactical sequencing rather than
in single-step translation.

Raw logs:
\texttt{logs/rl\_compare/our\_deepseek\_results\_runs\_v4flash\_n30.csv},
summary:
\texttt{logs/rl\_compare/our\_deepseek\_summary\_v4flash\_n30.csv}.

\subsection{7. Discussion}\label{discussion}

The main lesson is that LLMs do not need to operate at the same level as
traditional game-control agents. They are poorly matched to
high-frequency per-tick control and fragile numerical battlefield
prediction. They are better matched to translating natural language into
structured, auditable, task-level commands.

The co-pilot framing also changes what the system should and should not
do. A traditional autonomous agent may need to estimate whether an
attack is likely to succeed. This system should not. The player sees the
game, makes the strategic judgment, and decides what they want. The AI
co-pilot should execute the declared tactical intent faithfully.

The project also suggests a plausible path toward local models. The
output space is structured MCP JSON over a small tool vocabulary. This
is a better distillation target than open-ended strategic reasoning. A
large model can generate high-quality demonstrations, and a smaller
model can later be trained to produce valid MCP calls under lower
latency and cost.

\subsection{8. Limitations}\label{limitations}

This is an early prototype and should be presented as such.

\textbf{No formal user study.} All evaluations in §6 are LLM-driven from
scripted natural-language commands, not from human players unfamiliar
with the system. We do not measure player learning curve, subjective
ease-of-use, recovery from misinterpretation, or per-task wallclock
versus manual mouse-and-keyboard play. Claims of ``plain natural
language'' are supported by the breadth of capability tests (T1--T10,
spanning referent resolution, kind/state splitting, mid-flight
re-commanding, partial cancel, conditional trigger, route constraint,
formation, time-sequencing, failure recovery) and by a live free-form
LLM session, but a controlled human-subjects study remains future work
(see §9).

\textbf{Small sample sizes.} The repeated DeepSeek trials are N=5 per
scenario; the cross-paradigm cost reference is N=3 full games. These are
sufficient to characterise central tendency and basic variance but not
for strong statistical claims. We report mean ± standard deviation and
the per-run table to allow readers to assess noise directly.

\textbf{Workload asymmetry in the cost reference.} The OpenRA-RL cost
reference in §6.6 plays a complete game (economy + combat + outcome),
whereas the openra\_mcp scenarios in §6.5 execute isolated tactical
commands on a pre-trained roster. This is by design --- player-owned
economy is a paradigm constraint, not an evaluation oversight --- but
makes the two cost numbers not directly comparable as a win-rate or
end-to-end-efficiency benchmark.

\textbf{Scenario coverage.} Tactical scenarios use prepared OpenRA
setups rather than full competitive matches; sandbox conditions and
\texttt{/instantbuild} cheats accelerate roster preparation. Evaluation
has been conducted on the Soviet faction; Allied units
(1tnk/2tnk/medi/mech) have not been systematically retested
post-ablation.

\textbf{Engine-side limitations.} The squad-overlap interaction
documented in §6.5 --- when the LLM dispatches a new squad without first
cancelling the holding squad --- is a real system limitation, not a
translation failure. Improving the engine to either reassign units on
overlap or surface the conflict to the LLM is engineering work.

\textbf{Information boundary.} Player-owned-information and
player-owned- economy boundaries are currently enforced by protocol and
tool surface design, not by a complete formal security mechanism.

\textbf{Local small model not yet trained.} The future-work distillation
path is unrealised in this paper.

\textbf{Telemetry mixes sources.} Development-session telemetry (§6.4)
includes debugging and authoring sessions and should be read as
prototype observability, not benchmark data.

\subsection{9. Future Work}\label{future-work}

Future work should proceed in four directions. First, run repeated clean
evaluations across maps, factions, and enemy conditions. Second, compare
the co-pilot against manual control using click count, completion time,
and error rate. Third, extract supervised training pairs from LLM
demonstrations and fine-tune a small local model for low-latency MCP
command generation. Fourth, strengthen the boundary between player-owned
economic/information authority and AI-owned tactical translation.

\subsection{10. Conclusion}\label{conclusion}

\texttt{openra\_mcp} demonstrates a practical early version of
natural-language tactical control for RTS games. The system does not try
to make the LLM a complete player. Instead, it keeps the human in charge
of strategy, economy, and battlefield judgment, while the LLM translates
tactical intent into auditable MCP calls over a small set of squad
execution primitives. The current prototype evidence is preliminary but
concrete: complex OpenRA maneuvers such as splitting, pincer movement,
route constraints, time sequencing, and failure recovery can be
expressed as player language and executed in the game.

The broader implication is that language models may be most useful in
games when they are placed at the right level of abstraction. They need
not drive every unit every tick. For RTS co-piloting, the more promising
role is to turn human tactical intent into structured, executable
control while the engine handles local behavior. This framing also
creates a clear path toward small local models: the learning target is
not full autonomous strategy, but valid low-latency MCP command
generation from player intent and compact game state.

\subsection{Data and Code
Availability}\label{data-and-code-availability}

This preprint and its supporting artifacts are archived on Zenodo. The
current version (v2) DOI is \texttt{10.5281/zenodo.20393182}
(https://zenodo.org/records/20393182). The concept DOI
\texttt{10.5281/zenodo.20377061} always resolves to the latest version.
The release package (\texttt{000\_UPLOAD\_TO\_ZENODO\_openra\_mcp.zip})
contains the manuscript (\texttt{.pdf} / \texttt{.tex} / \texttt{.docx}
/ \texttt{.md}), the BibTeX references, a source-code snapshot
(\texttt{mcp\_server/}, \texttt{trait\_src/}), per-run evaluation CSVs
(\texttt{logs/v2\_post\_ablation.csv},
\texttt{logs/v2\_retry\_4fails.csv}, \texttt{logs/v2\_t7\_retry2.csv},
\texttt{logs/rl\_compare/our\_deepseek\_results\_runs\_n30.csv},
\texttt{logs/rl\_compare/our\_deepseek\_results\_runs\_n100.csv},
\texttt{logs/rl\_compare/rl\_full\_game\_n3\_summary.csv},
\texttt{logs/rl\_compare/rl\_scripted\_n3\_summary.csv}),
telemetry-derived metrics (\texttt{logs/paper\_metrics.json}), and
supplementary notes. Demonstration videos are listed in
\texttt{papers/SUPPLEMENTARY\_MATERIALS.md} and are intentionally not
included in the core release zip due to their size.

The evaluation evidence used in this preprint is drawn from:

\begin{itemize}
\tightlist
\item
  \texttt{mcp\_server/experiments/scenarios\_v2.py}
\item
  \texttt{logs/v2\_post\_ablation.csv}
\item
  \texttt{logs/v2\_retry\_4fails.csv}
\item
  \texttt{logs/v2\_t7\_retry2.csv}
\item
  \texttt{logs/baseline\_pre\_ablation.md}
\item
  \texttt{logs/live\_llm\_demo/demo\_01.mp4}
\item
  \texttt{logs/v2\_videos/*.mp4}
\item
  \texttt{logs/paper\_metrics.json}
\end{itemize}

The first public release should include the manuscript, the relevant CSV
files, the code snapshot, and either the videos themselves or stable
links to the video supplement.

\subsection{Ethics and Boundary
Statement}\label{ethics-and-boundary-statement}

This project is a game-control prototype and does not involve
human-subjects data collection in its current evaluation. The main
ethical issue is not player privacy but control delegation. The system
is therefore designed around an explicit boundary: the player owns
strategy, economy, and battlefield judgment, while the LLM translates
tactical intent into executable actions. The paper does not claim that
the system should make strategic decisions for players or that it can
predict whether an engagement is favorable.

\subsection{AI Assistance Disclosure}\label{ai-assistance-disclosure}

This manuscript and its supporting metadata were drafted with AI
assistance. The author provided the project, design intent,
implementation evidence, and publication goal. AI assistance was used to
organize the argument, draft prose, prepare tables from local logs, and
assemble release metadata. The author is responsible for verifying the
final claims, references, and release artifacts.

\subsection{References}\label{references}

Ahn, D., Kim, S., \& Choi, J. (2025). \emph{Society of mind meets
real-time strategy: A hierarchical multi-agent framework for strategic
reasoning}. arXiv:2508.06042. https://arxiv.org/abs/2508.06042

Amershi, S., Weld, D., Vorvoreanu, M., Fourney, A., Nushi, B.,
Collisson, P., Suh, J., Iqbal, S., Bennett, P. N., Inkpen, K., Teevan,
J., Kikin-Gil, R., \& Horvitz, E. (2019). Guidelines for human-AI
interaction. In \emph{Proceedings of the 2019 CHI Conference on Human
Factors in Computing Systems} (pp.~1-13).
https://doi.org/10.1145/3290605.3300233

Anne, T., Syrkis, N., Elhosni, M., Turati, F., Legendre, F., Jaquier,
A., \& Risi, S. (2025). \emph{Harnessing language for coordination: A
framework and benchmark for LLM-driven multi-agent control}.
arXiv:2412.11761. https://arxiv.org/abs/2412.11761

Anthropic. (2024). \emph{Introducing the Model Context Protocol}.
https://www.anthropic.com/research/model-context-protocol

Horvitz, E. (1999). Principles of mixed-initiative user interfaces. In
\emph{Proceedings of the SIGCHI Conference on Human Factors in Computing
Systems} (pp.~159-166). https://doi.org/10.1145/302979.303030

HOPPINZQ. (2025). \emph{hoppinai-mcp-red95: Conversational Red Alert via
MCP}. Open-source software, first commit 2025-10-29.
https://github.com/HOPPINZQ/hoppinai-mcp-red95

Iannoli, A., Gigli, L., Sciullo, L., Trotta, A., \& Di Felice, M.
(2026). \emph{Say the mission, execute the swarm: Agent-enhanced LLM
reasoning in the Web-of-Drones}. arXiv:2605.03788.
https://arxiv.org/abs/2605.03788

Liu, J., Yu, C., Gao, J., Xie, Y., Liao, Q., Wu, Y., \& Wang, Y. (2023).
\emph{LLM-powered hierarchical language agent for real-time human-AI
coordination}. arXiv:2312.15224. https://arxiv.org/abs/2312.15224

Ma, W., Mi, Q., Zeng, Y., Yan, X., Wu, Y., Lin, R., Zhang, H., \& Wang,
J. (2024). \emph{Large language models play StarCraft II: Benchmarks and
a chain of summarization approach}. arXiv:2312.11865.
https://arxiv.org/abs/2312.11865

OpenRA Project. (n.d.). \emph{About OpenRA}.
https://www.openra.net/about/

OpenRA-RL Contributors. (2026). \emph{OpenRA-RL: Command AI to play Red
Alert}. https://openra-rl.dev/

Ramos-Silva, J. N., \& Burke, P. J. (2026). \emph{A universal large
language model: Drone command and control interface}. arXiv:2601.15486.
https://arxiv.org/abs/2601.15486

Shao, X., Jiang, W., Zuo, F., \& Liu, M. (2024). \emph{SwarmBrain:
Embodied agent for real-time strategy game StarCraft II via large
language models}. arXiv:2401.17749. https://arxiv.org/abs/2401.17749

Vinyals, O., Babuschkin, I., Czarnecki, W. M., Mathieu, M., Dudzik, A.,
Chung, J., Choi, D. H., Powell, R., Ewalds, T., Georgiev, P., Oh, J.,
Horgan, D., Kroiss, M., Danihelka, I., Huang, A., Sifre, L., Cai, T.,
Agapiou, J. P., Jaderberg, M., \& Silver, D. (2019). Grandmaster level
in StarCraft II using multi-agent reinforcement learning. \emph{Nature,
575}(7782), 350-354. https://doi.org/10.1038/s41586-019-1724-z

Wang, G., Xie, Y., Jiang, Y., Mandlekar, A., Xiao, C., Zhu, Y., Fan, L.,
\& Anandkumar, A. (2023). \emph{Voyager: An open-ended embodied agent
with large language models}. arXiv:2305.16291.
https://arxiv.org/abs/2305.16291

\end{document}
